\scanpage{248}
\S~1 Un théorème de géométrie analytique.

Rappelons d'abord deux résultats connus

\uline{Lemme} 1.1. \unsure{Soient|Soit} $X$ un espace \deleted{\unsure{para|pour}} topologique, $Y$
une partie fermée de $X$ ayant un syst. \unsure{fondamental|fond.mental}
de voisinages paracompacts, $F$ un faisceau abélien
sur $X$, $F|Y$ le faisceau induit sur $Y$. Alors
l'hom. can. \deleted{\unsure{est}} suivant est un isom :
\[
\underset{U}{\underrightarrow{\lim}}H^*(U,F)\longrightarrow\unsure{-|\vphantom{-}}H^*(Y,F|Y)
\]
[où dans le premier \unsure{mb|mt|membre}, la \uline{lim} est sur
l'ensemble des voisinages de $Y$ \unsure{ds|de} $X$]

\uline{Corollaire} 1.2 Soit $f:X\to Y$ une application
continue \deleted{propre}\inserted{fermée}, avec $Y$ loc compact,
\deleted{$F$ un} $y\in Y$, $X_0=f^{-1}(\{y\})$, $F$ un faisceau
abélien sur $X$, $F|X_0$ le faisceau induit
sur $X_0$. Alors \unsure{on a} un isom
canonique
\[
R^*f_{*}(F)_y\simeq H^*(X_{\unsure{0|y}},F|X_y)
\]
% archive sheet number: 247
